\subsubsection{有限 RH：\texorpdfstring{$s$}{s}-变量 $\zeta$ 的临界带极点与 primitive 的平方根振荡}\label{app:real-input-40-finite-rh}
把 $z$-变量换为 Dirichlet 坐标
\[
z=\lambda^{-s},\qquad \widehat{\zeta}(s):=\zeta_M(\lambda^{-s})=\frac{1}{\det(I-\lambda^{-s}M)}.
\]
则 $\widehat{\zeta}(s)$ 的极点由 $M$ 的非零特征值 $\mu$ 决定：当且仅当
\[
\lambda^{-s}\mu=1
\]
时出现极点（对每个 $\mu$ 及其对数分支给出一列极点）。

\paragraph{谱—极点—误差三角形：有限维 \texorpdfstring{$\zeta$}{zeta} 的通用字典}
对任意有限维核（矩阵）$M$，Dirichlet 化后的 $\widehat{\zeta}(s)$ 把“谱模长”逐点翻译为“临界带极点的位置”，并进一步把“非平凡谱半径”翻译为 primitive 计数的误差尺度。
该字典是有限维模型中最接近经典 \(\zeta\)/\(L\) 语言的部分：它把“零点/极点—显式公式—误差项”的分析数论三角形压缩为一个纯线性代数三角形。

\begin{proposition}[谱 \texorpdfstring{$\leftrightarrow$}{<->} 极点：临界带坐标与 toy-RH 判别]\label{prop:finite-rh-triangle-dict}
设 $M$ 为有限维矩阵，$\lambda=\rho(M)>1$，并定义
\[
\widehat{\zeta}(s)=\zeta_M(\lambda^{-s})=\frac{1}{\det(I-\lambda^{-s}M)}.
\]
则对任意非零特征值 $\mu\neq 0$ 与任意 $k\in\ZZ$，$\widehat{\zeta}(s)$ 的极点由
\[
\lambda^{-s}\mu=1
\]
给出，等价于
\[
s=\frac{\log\mu+2\pi i k}{\log\lambda},
\]
其中 $\log\mu$ 取任意复对数分支。于是极点实部满足
\[
\Re(s)=\frac{\log|\mu|}{\log\lambda}.
\]
因此在开带 $0<\Re(s)<1$ 中出现非平凡极点当且仅当存在特征值满足 $1<|\mu|<\lambda$。

进一步，若把“有限维 toy-RH”定义为：开带 $0<\Re(s)<1$ 中的全部非平凡极点都落在 $\Re(s)=\tfrac12$，则该条件等价于
\[
(\forall\,\mu\ \text{满足}\ 1<|\mu|<\lambda)\qquad |\mu|=\sqrt{\lambda}.
\]
更弱地，若满足 Ramanujan 半径界
\[
\max_{\mu\neq\lambda}|\mu|\ \le\ \sqrt{\lambda},
\]
则右半开带 $\tfrac12<\Re(s)<1$ 内无极点；若存在 $|\mu|=\sqrt{\lambda}$，则对应极点列恰落在 $\Re(s)=\tfrac12$ 上。
\end{proposition}

\begin{remark}[本核的达界位置：唯一临界模长 \texorpdfstring{$|\mu|=\sqrt{\lambda}$}{|mu|=sqrt(lambda)}]\label{rem:finite-rh-40-ramanujan}
在真实输入 40 状态核上，除 Perron 根 $\lambda=\varphi^2$ 外，唯一满足 $|\mu|=\sqrt{\lambda}$ 的特征值为 $-\sqrt{\lambda}=-\varphi$（命题 \ref{prop:real-input-40-essential-20} 与此处命题 \ref{prop:finite-rh-40} 的谱分解一致）。
因此该核达到 toy-RH 的“单频临界振荡”情形：临界带中的全部非平凡极点由 $-\sqrt{\lambda}$ 单独产生。
\end{remark}

\paragraph{相位增殖：保持 RH-sharp 的循环提升}
经典 \(\zeta(s)\) 的困难部分不仅是“实部落在 \(\tfrac12\)”这一条临界线，而是临界线上存在大量、且虚部结构复杂的零点族。
在有限维核模型中，可以在\textbf{保持 Ramanujan/RH-sharp 半径界}的同时，把“单一相位”的临界极点列增殖为任意细的相位梳齿；这为后续把 toy-RH 推向 toy-GRH（通过群扩张/Artin 分解）提供了一个最小可审计机制。

\begin{proposition}[循环相位提升：谱相位细分且保持 Ramanujan 半径界]\label{prop:finite-rh-phase-lift}
设 $M$ 为有限维矩阵，$\lambda=\rho(M)>1$。取整数 $q\ge 2$，令 $P_q$ 为 $q$-循环置换矩阵（$P_q e_j=e_{j+1}$，指标按 $q$ 取模），并定义循环提升矩阵
\[
M^{\langle q\rangle}:=M\otimes P_q.
\]
则：
\begin{enumerate}
  \item \textbf{（谱相位细分）} 设 $\omega_q=e^{2\pi i/q}$。有谱恒等式
  \[
  \mathrm{spec}(M^{\langle q\rangle})
  =\{\mu\,\omega_q^k:\ \mu\in \mathrm{spec}(M),\ k=0,1,\dots,q-1\},
  \]
  按代数重数计。
  \item \textbf{（谱半径保持）} $\rho(M^{\langle q\rangle})=\rho(M)=\lambda$。
  \item \textbf{（Ramanujan 界保持）} 若 $M$ 满足 $\max_{\mu\neq \lambda}|\mu|\le \sqrt{\lambda}$，则 $M^{\langle q\rangle}$ 亦满足同一半径界。
  \item \textbf{（临界线极点梳齿）} 令
  \(
  \widehat{\zeta}^{\langle q\rangle}(s):=\det(I-\lambda^{-s}M^{\langle q\rangle})^{-1}
  \)。
  若 $\mu=\sqrt{\lambda}e^{i\theta}$ 是 $M$ 的某个达界特征值，则它在 $M^{\langle q\rangle}$ 上生成 $q$ 条互相平移的临界线极点列：
  \[
  s=\frac12+\frac{i}{\log\lambda}\Bigl(\theta+2\pi\Bigl(\frac{k}{q}+m\Bigr)\Bigr),
  \qquad k\in\{0,\dots,q-1\},\ m\in\ZZ.
  \]
\end{enumerate}
当 $q\to\infty$ 时，上述极点在每个基本周期窗 $[0,2\pi/\log\lambda)$ 内的虚部位置变得任意稠密。
\end{proposition}

\begin{proof}
置换矩阵 $P_q$ 可被离散傅里叶变换对角化，其特征值为 $\omega_q^k$（$0\le k\le q-1$）。
对任意 $M$ 的特征对 $(\mu,v)$ 与任意 $P_q$ 的特征对 $(\omega_q^k,w_k)$，张量向量 $v\otimes w_k$ 满足
\(
(M\otimes P_q)(v\otimes w_k)=(Mv)\otimes(P_qw_k)=\mu\,\omega_q^k (v\otimes w_k)
\)，从而得到谱相位细分公式；反向包含由维数计数或 Kronecker 谱性质得到。
谱半径与 Ramanujan 界的保持性来自 $|\omega_q^k|=1$。
最后极点公式由命题 \ref{prop:finite-rh-triangle-dict} 对 $M^{\langle q\rangle}$ 逐个特征值应用并取达界模长 $|\mu|=\sqrt{\lambda}$ 得到。证毕。
\end{proof}

\begin{remark}[应用到真实输入 $40$ 状态核：单频 \texorpdfstring{$\pi$}{pi} 振荡的细分]\label{rem:finite-rh-40-phase-lift}
对真实输入 $40$ 状态核，唯一达界特征值为 $-\sqrt{\lambda}=-\varphi=\sqrt{\lambda}e^{i\pi}$（注 \ref{rem:finite-rh-40-ramanujan}）。
因此循环相位提升把原来的单一临界列
\(
s=\tfrac12+\frac{(2m+1)\pi i}{\log\lambda}
\)
细分为
\[
s=\frac12+\frac{i}{\log\lambda}\Bigl(\pi+2\pi\Bigl(\frac{k}{q}+m\Bigr)\Bigr),
\qquad k=0,\dots,q-1,\ m\in\ZZ,
\]
在保持 RH-sharp 的同时把临界线上的虚部结构从“单栅栏”提升为“$q$-梳齿”。
\end{remark}

\begin{proposition}[循环提升的 Artin 因子化（阿贝尔角色通道的显式分解）]\label{prop:finite-rh-phase-lift-artin}
沿用命题 \ref{prop:finite-rh-phase-lift} 的记号。令
\(
\omega_q=e^{2\pi i/q}
\)，并把循环群 \(G=\ZZ/q\ZZ\) 的角色记为 \(\chi_k(j)=\omega_q^{kj}\)（\(k\in\{0,\dots,q-1\}\)）。
则对任意复参数 \(z\) 有行列式分解
\[
\det\!\bigl(I-zM^{\langle q\rangle}\bigr)
=\prod_{k=0}^{q-1}\det\!\bigl(I-z\,\omega_q^k M\bigr).
\]
等价地，循环提升的 \(\zeta\)-函数满足
\[
\zeta_{M^{\langle q\rangle}}(z)
=\prod_{k=0}^{q-1}\zeta_{M}(z\omega_q^k),
\qquad
\zeta_M(z):=\det(I-zM)^{-1}.
\]
因此“相位增殖”不仅是谱点集合的旋转细分，更是一个可审计的\textbf{阿贝尔 Artin 分解}：每个角色通道 \(\chi_k\) 对应一个“扭曲因子” \(\det(I-z\,\omega_q^k M)\)。
\end{proposition}

\begin{proof}
取离散傅里叶矩阵 \(F\) 使得 \(F^{-1}P_qF=\mathrm{diag}(\omega_q^0,\dots,\omega_q^{q-1})\)。
则
\[
I-z(M\otimes P_q)
\sim
I-z\bigl(M\otimes \mathrm{diag}(\omega_q^0,\dots,\omega_q^{q-1})\bigr)
\]
在相似意义下等价于一个按 \(k\) 分块的块对角矩阵，其第 \(k\) 个块为 \(I-z\,\omega_q^k M\)。对块对角矩阵取行列式并相乘即得结论。证毕。
\end{proof}

\begin{proposition}[临界带中的非平凡极点都在 \texorpdfstring{$\Re(s)=\tfrac12$}{Re(s)=1/2}]\label{prop:finite-rh-40}
设 $\lambda=\rho(M)=\varphi^2$。则在开带 $0<\Re(s)<1$ 中，$\widehat{\zeta}(s)$ 的全部极点均严格落在直线
\[
\Re(s)=\frac12.
\]
更具体地，唯一落在 $0<\Re(s)<1$ 的极点列来自特征值 $-\varphi=-\sqrt{\lambda}$，其位置为
\[
s=\frac12+\frac{(2k+1)\pi i}{\log\lambda},\qquad k\in\ZZ.
\]
\end{proposition}

\begin{proof}
由上文的谱分解，$M$ 的非零特征值集合为
\[
\{\lambda,\ -\sqrt{\lambda},\ 1^{(2)},\ (-1)^{(3)},\ \lambda^{-1},\ \lambda^{-1/2}\},
\]
其余 $31$ 个特征值为 $0$。对每个非零特征值 $\mu$，$\det(I-\lambda^{-s}M)=0$ 等价于 $\lambda^{-s}\mu=1$。

当 $\mu=\lambda$ 时得到主极点列 $s=1+2\pi i k/\log\lambda$，不落在开带 $0<\Re(s)<1$。
当 $\mu=-\sqrt{\lambda}$ 时，取复对数分支 $\log(-\sqrt{\lambda})=\tfrac12\log\lambda+i(2k+1)\pi$，
于是得到
\(
s=\tfrac12+\frac{(2k+1)\pi i}{\log\lambda}
\)，其满足 $0<\Re(s)<1$。
其余特征值满足 $|\mu|\le 1$，对应的极点满足 $\Re(s)\le 0$（落在临界带左侧边界或更左），因此在开带 $0<\Re(s)<1$ 内不产生极点。
\end{proof}

\begin{remark}[Dirichlet 变量与“临界线”表述]
令 $r:=\lambda^{1-s}$，则因子 $r+\varphi$ 对应的奇点满足
\[
r=-\varphi=\varphi e^{i(\pi+2\pi k)}.
\]
由 $\lambda=\varphi^2$ 得
\[
s=\frac12-\frac{(2k+1)\pi i}{\log\lambda},
\]
即在 $\Re(s)>0$ 内的非实奇点全部落在 $\Re(s)=1/2$ 上。
\end{remark}

\begin{proposition}[RH-sharp 的奇偶消项：Möbius 反演的一般机制]\label{prop:finite-rh-parity-general}
设 \(\lambda>1\)，并给定一列实数 \(a_n\) 满足
\[
a_n=\lambda^n+\epsilon^n\lambda^{n/2}+O(\lambda^{n/3}),
\qquad \epsilon\in\{\pm 1\}.
\]
定义 primitive 轨道数（Witt/Möbius 反演）
\[
p_n:=\frac{1}{n}\sum_{d\mid n}\mu(d)\,a_{n/d}.
\]
则当 \(n\to\infty\) 时有：
\begin{enumerate}
  \item 若 \(n\) 为奇数，则
  \[
  p_n=\frac{\lambda^n}{n}+\frac{\epsilon^n\lambda^{n/2}}{n}+O\!\left(\frac{\lambda^{n/3}}{n}\right).
  \]
  \item 若 \(n\) 为偶数，则 \(d=1,2\) 两项在 \(\lambda^{n/2}\) 级别发生精确抵消，从而
  \[
  p_n=\frac{\lambda^n}{n}+O\!\left(\frac{\lambda^{n/3}}{n}\right).
  \]
\end{enumerate}
\end{proposition}
\begin{proof}
由定义
\[
np_n=\sum_{d\mid n}\mu(d)\,a_{n/d}.
\]
若 \(n\) 为奇数，则对所有 \(d\mid n\) 且 \(d\ge 2\) 有 \(d\ge 3\)，因此
\(
a_{n/d}=O(\lambda^{n/3})
\)。
于是
\[
np_n=a_n+\sum_{\substack{d\mid n\\ d\ge 3}}\mu(d)\,a_{n/d}
=\lambda^n+\epsilon^n\lambda^{n/2}+O(\lambda^{n/3}),
\]
得第一式。
\par
若 \(n\) 为偶数，则把 \(d=1,2\) 两项分离：
\[
np_n=a_n-a_{n/2}+\sum_{\substack{d\mid n\\ d\ge 3}}\mu(d)\,a_{n/d}.
\]
由假设，因 \(n\) 为偶数而 \(\epsilon^n=1\)，故
\[
a_n-a_{n/2}
\;=\;
(\lambda^n+\lambda^{n/2}+O(\lambda^{n/3}))-(\lambda^{n/2}+O(\lambda^{n/4}))
\;=\;
\lambda^n+O(\lambda^{n/3}).
\]
并且余项仍为 \(O(\lambda^{n/3})\)，合并得第二式。证毕。
\end{proof}

\begin{corollary}[显式公式形状：\texorpdfstring{$a_n$}{a\_n} 的平方根振荡与 primitive 的奇偶消项]\label{cor:finite-rh-explicit-formula}
记 $a_n=\Tr(M^n)$。则
\[
a_n=\lambda^n+(-1)^n\lambda^{n/2}+O(1),
\]
并且 primitive 轨道数满足更精细的奇偶结构：
\[
p_n=\frac{\lambda^n}{n}-\frac{\lambda^{n/2}}{n}+O\!\left(\frac{\lambda^{n/3}}{n}\right)\qquad(n\ \text{为奇数}),
\]
\[
p_n=\frac{\lambda^n}{n}+O\!\left(\frac{\lambda^{n/3}}{n}\right)\qquad(n\ \text{为偶数}),
\]
其中偶数情形的 $\lambda^{n/2}$ 级项被 Möbius 反演中的 $d=1$ 与 $d=2$ 两项\textbf{精确抵消}。
\[
\boxed{
\;p_n=\frac{\lambda^n}{n}+\frac{\bigl((-1)^n-\mathbf{1}_{2\mid n}\bigr)\lambda^{n/2}}{n}
+O\!\left(\frac{\lambda^{n/3}}{n}\right)\; }
\]
并且具有可核验的“奇偶共振 signature”：
\[
\lim_{\substack{n\to\infty\\ n\ \mathrm{odd}}}\ \frac{n}{\lambda^{n/2}}\left(p_n-\frac{\lambda^n}{n}\right)=-1,
\qquad
\lim_{\substack{n\to\infty\\ n\ \mathrm{even}}}\ \frac{n}{\lambda^{n/2}}\left(p_n-\frac{\lambda^n}{n}\right)=0.
\]
\end{corollary}

\begin{remark}[偶长层的“超强化”主次项]\label{rem:finite-rh-even-strengthen}
在偶长层，$\lambda^{n/2}$ 项由 $d=1$ 与 $d=2$ 精确抵消。进一步分情形可得更强的主次项：
\[
p_n=\frac{\lambda^n}{n}-\frac{\lambda^{n/3}}{n}+O\!\left(\frac{\lambda^{n/4}}{n}\right),
\qquad n\ \text{偶且}\ 3\mid n,
\]
\[
p_n=\frac{\lambda^n}{n}+(-1)^{n/2+1}\frac{\lambda^{n/4}}{n}
O\!\left(\frac{\lambda^{n/5}}{n}\right),
\qquad n\ \text{偶且}\ 3\nmid n.
\]
其中第一式的 $-\lambda^{n/3}$ 来自 $d=3$ 项（$\mu(3)=-1$），第二式的 $\lambda^{n/4}$ 来自 $d=2$ 项中的次级振荡
$(-\varphi)^{n/2}=(-1)^{n/2}\lambda^{n/4}$；而 $\mu(4)=0$ 使 $d=4$ 不贡献主次级项。
\end{remark}

\begin{remark}[次级项由“谱信息 $\times$ 因子结构”共同锁定]\label{rem:finite-rh-factor-controls-secondary}
在本核上，唯一满足 $|\mu|=\sqrt{\lambda}$ 的非平凡特征值为 $-\sqrt{\lambda}$，因此 primitive 侧的“平方根级误差”与“偶长层相消”都由 Möbius 反演的 $d=1,2$ 两项主导（推论 \ref{cor:finite-rh-explicit-formula}）。
一旦该层级被抵消，新的主次级项不再只由谱决定，还会显式依赖 $n$ 的因子结构：对 $n$ 的任意奇素因子 $p\mid n$，Möbius 反演中的 $d=p$ 项会贡献一个量级为 $\lambda^{n/p}/n$ 的修正（其符号由 $\mu(p)=-1$ 决定）。
\par
更紧凑地，定义
\[
p_{\min}(n):=\min\{p:\ p\ \text{为奇素数且}\ p\mid n\},
\qquad(\text{若无奇素因子则令 }p_{\min}(n)=+\infty).
\]
则对偶数 $n$，在 $d=1,2$ 的 $\lambda^{n/2}$ 级精确相消之后，剩余次级项的指数尺度由
\[
\max\Bigl(\frac14,\ \frac{1}{p_{\min}(n)}\Bigr)
\]
决定：当 $3\mid n$ 时由 $d=3$ 项给出 $\lambda^{n/3}/n$，当 $3\nmid n$ 时由 $d=2$ 项中的次级振荡给出 $\lambda^{n/4}/n$，而更大奇素因子只会产生更低阶的 $\lambda^{n/p}$ 修正。
\end{remark}

\begin{corollary}[RH-型误差的锐性（长度与 $x$ 口径）]\label{cor:finite-rh-prime-orbit-osc}
令
\[
\Pi(N):=\sum_{n\le N}p_n,\qquad x:=\lambda^N.
\]
由推论 \ref{cor:finite-rh-explicit-formula} 得
\[
\Pi(N)=\frac{\lambda}{\lambda-1}\cdot\frac{\lambda^N}{N}+O\!\left(\frac{\lambda^{N/2}}{N}\right),
\]
且由于 $(-\varphi)$ 的实际存在，上式中的平方根量级不可改进到 $\lambda^{(1/2-\varepsilon)N}$。等价地
\[
\pi(x)=\frac{\lambda\log\lambda}{\lambda-1}\cdot\frac{x}{\log x}
\;+\Omega_{\pm}\!\left(\frac{x^{1/2}}{\log x}\right),
\]
并且振荡频率由 $\Im(s)=\frac{(2k+1)\pi}{\log\lambda}$ 给出（周期 $2\pi/\log\lambda$）。
\end{corollary}

\begin{proof}
由上文的闭式
\(
a_n=\varphi^{2n}+(-\varphi)^n+2+3(-1)^n+\varphi^{-2n}+\varphi^{-n}
\)
可得
\(
a_n=\lambda^n+(-1)^n\lambda^{n/2}+O(1)
\)。

由 Möbius 反演 $p_n=\frac1n\sum_{d\mid n}\mu(d)a_{n/d}$：
若 $n$ 为奇数，则所有 $d\ge 2$ 满足 $d\ge 3$，故 $a_{n/d}=O(\lambda^{n/3})$。于是
\[
np_n=a_n+\sum_{\substack{d\mid n\\ d\ge 3}}\mu(d)a_{n/d}
=\lambda^n-\lambda^{n/2}+O(\lambda^{n/3}),
\]
得到第一式。

若 $n$ 为偶数，则写出 $d=1,2$ 两项：
\[
np_n=a_n-\;a_{n/2}+\sum_{\substack{d\mid n\\ d\ge 3}}\mu(d)a_{n/d}.
\]
此时 $a_n=\lambda^n+\lambda^{n/2}+O(1)$（因为 $(-1)^n=1$），而 $a_{n/2}=\lambda^{n/2}+O(\lambda^{n/4})$，故
\[
a_n-a_{n/2}=\lambda^n+O(\lambda^{n/4}),
\]
并且余项由 $d\ge 3$ 给出 $O(\lambda^{n/3})$，从而得到第二式。这里 $\lambda^{n/2}$ 的抵消正是 $a_n$ 的平方根项（$d=1$）与 $a_{n/2}$ 的主项（$d=2$）之间的精确抵消。
\end{proof}

\begin{theorem}[primitive--Dirichlet 级数的自然边界：\texorpdfstring{$\Re(s)=0$}{Re(s)=0}]\label{thm:real-input-40-primedirichlet-dense-branch}
令
\[
\mathcal{P}(s):=P(\lambda^{-s})=\sum_{n\ge 1}p_n\,\lambda^{-ns}\qquad(\Re(s)>1),
\]
其中 $P(z)=\sum_{n\ge 1}p_n z^n$ 为 primitive 生成函数，$\lambda=\rho(M)=\varphi^2$，并定义
\[
\widehat{\zeta}(s):=\zeta_M(\lambda^{-s}).
\]
则：
\begin{enumerate}
  \item \textbf{（解析延拓到右半平面）} 对任意 $\sigma_0>0$，级数
  \[
  \sum_{k\ge 1}\frac{\mu(k)}{k}\log \widehat{\zeta}(ks)
  \]
  在集合
  \(
  \{s:\Re(s)\ge \sigma_0\}
  \)
  的任意紧子集上（避开下述极点列）一致收敛，从而定义出 $\mathcal{P}(s)$ 在开半平面 $\Re(s)>0$ 上的一个解析延拓（其在 $\Re(s)>1$ 与原 Dirichlet 级数一致）。
  \item \textbf{（稠密奇点列与自然边界）} 对每个奇素数 $p$ 与任意 $m\in\ZZ$，点
  \[
  s_{p,m}:=\frac{1}{p}+\frac{2\pi i m}{p\log\lambda}
  \]
  是该解析延拓的一个对数型分支奇点。并且集合 $\{s_{p,m}\}_{p\ \mathrm{odd\ prime},\,m\in\ZZ}$ 在直线 $\Re(s)=0$ 上稠密。
\end{enumerate}
因此 \(\Re(s)=0\) 是 \(\mathcal{P}(s)\) 的\textbf{自然边界}：\(\mathcal{P}(s)\) 不可能在任何虚轴点的邻域内作解析延拓穿过该直线。
\end{theorem}

\begin{proof}
\emph{第 1 点（收敛与解析性）}：
当 $|z|$ 足够小时有 $\log\zeta_M(z)=O(z)$（因为 $\zeta_M(0)=1$ 且 $\zeta_M$ 在 $0$ 的邻域解析），因此对任意 $\sigma_0>0$，当 $\Re(s)\ge \sigma_0$ 时
\(
|\log\widehat{\zeta}(ks)|=|\log\zeta_M(\lambda^{-ks})|=O(\lambda^{-k\sigma_0})
\)
随 $k$ 指数衰减。于是 $\sum_{k\ge 1}|\mu(k)|\,|\log\widehat{\zeta}(ks)|/k$ 在任意避开极点列的紧集上一致收敛；逐项可导与一致收敛给出其和函数在 $\Re(s)>0$（去掉极点列）上解析，从而得到 $\mathcal{P}(s)$ 的解析延拓。

\emph{第 2 点（显式奇点与稠密性）}：
由本核的行列式分解
\[
\det(I-zM)=(1-z)^2(1+z)^3(1-3z+z^2)(1+z-z^2),
\]
可见 $z_\star=\lambda^{-1}$ 是 $\zeta_M(z)$ 的一个\textbf{简单极点}（对应因子 $1-3z+z^2$ 的简单零点），故 $\log\zeta_M(z)$ 在 $z=z_\star$ 具有对数型分支奇性。
当 $s=s_{p,m}$ 时有
\[
\lambda^{-ps}=\lambda^{-1}e^{-2\pi i m}=z_\star,
\]
因此 $\log\widehat{\zeta}(ps)=\log\zeta_M(\lambda^{-ps})$ 在 $s=s_{p,m}$ 处为对数型分支奇点，从而项
\(
\frac{\mu(p)}{p}\log\widehat{\zeta}(ps)
\)
在该点奇异。对任意 $k\neq p$，若 $\log\widehat{\zeta}(ks)$ 在 $s=s_{p,m}$ 奇异，则必有 $\lambda^{-ks_{p,m}}=z_\star$，等价于 $ks_{p,m}=1+\frac{2\pi i \ell}{\log\lambda}$（某个 $\ell\in\ZZ$）。比较实部得 $k/p=1$，故 $k=p$。因此除 $k=p$ 一项外，其余各项在 $s_{p,m}$ 的某邻域解析，且由第 1 点一致收敛可知其和亦解析，于是 $\mathcal{P}(s)$ 在 $s_{p,m}$ 必为对数型分支奇点。

最后证明稠密性：给定任意 $t\in\RR$ 与 $\varepsilon>0$，取足够大的奇素数 $p$，令
\(
m:=\left\lfloor \frac{t\,p\log\lambda}{2\pi}\right\rceil
\)
为最近整数，则
\[
\left|\Im(s_{p,m})-t\right|
=\left|\frac{2\pi m}{p\log\lambda}-t\right|
\le \frac{\pi}{p\log\lambda}<\varepsilon
\]
当 $p$ 足够大时成立，同时 $\Re(s_{p,m})=1/p<\varepsilon$。故 $s_{p,m}\to it$，从而 $\{s_{p,m}\}$ 在虚轴上稠密。由于任意穿过虚轴的邻域都包含这些奇点，故无法跨越解析延拓，\(\Re(s)=0\) 为自然边界。
\end{proof}

\begin{remark}[自然边界机制的普适化]\label{rem:primedirichlet-natural-boundary-universal}
定理 \ref{thm:real-input-40-primedirichlet-dense-branch} 的证明只用到了两点结构事实：
\begin{enumerate}
  \item $\widehat{\zeta}(s)=\zeta_M(\lambda^{-s})$ 在 $s=1$ 的极点来自 Perron 根 $\lambda$，从而 $\log\widehat{\zeta}(s)$ 在该点具有对数型分支奇性；
  \item prime-zeta 形式的恒等式
  \[
  \mathcal{P}(s)=\sum_{k\ge 1}\frac{\mu(k)}{k}\log\widehat{\zeta}(ks)
  \]
  把该分支奇性按比例 $s\mapsto s/k$ 复制到 $s_{p,m}=1/p+2\pi i m/(p\log\lambda)$（奇素数 $p$）处，并由 $k/p$ 的实部比较排除不同 $k$ 的抵消。
\end{enumerate}
因此该“稠密分支奇点 \(\Rightarrow\) \(\Re(s)=0\) 自然边界”的现象并不依赖本核的具体因式分解；
只要 $M$ 为 primitive 非负整数矩阵（保证 Perron 根简单）并以同样方式定义 primitive--Dirichlet 级数，完全同型的稠密奇点机制都会出现。
\end{remark}

\paragraph{与解析数论零密度技术的对齐（定位性备注）}
在经典 \(\zeta\) 理论中，“向 RH 推进”的一条可累积路线是改进 Dirichlet 多项式的大值控制，从而改进零密度界 \(N(\sigma,T)\) 并反馈到素数误差项的阈值。
例如 Guth--Maynard 给出新的大值估计并推出改进的零密度上界与短区间素数计数阈值（\(N(\sigma,T)\le T^{30(1-\sigma)/13+o(1)}\) 及相关推论，见 \cite{GuthMaynard2024LargeValues}）。

从“对齐口径”的角度，零密度技术的作用并不在于直接推出“全体零点都在临界线”，而在于把偏离临界线的零点集合压稀，从而把显式公式中的振荡误差项控制到更强的可积/可平均形式。
这对应于本文有限维原型中“最坏扭曲谱半径比 $\Lambda/\lambda_1$ 决定误差指数”的结构：当无法一步获得平方根界（核版 GRH），仍可通过控制“坏扭曲”的出现频率/规模来改进误差可控范围。

从本文的“有限维 \(\zeta\)”视角看，上述链条可被理解为一种连续谱版的“谱—极点—误差三角形”：零点（或极点）分布的稀疏性对应于显式公式中振荡项的可控叠加，而误差项的可达指数则由可用的大值界/矩母函数控制。有限维模型提供的优势在于：该三角形可被完全审计地压缩到有限谱数据（命题 \ref{prop:finite-rh-triangle-dict} 及推论 \ref{cor:finite-rh-explicit-formula}），从而为“把误差指数写成可检验谱不变量”提供了一个可计算原型。

\paragraph{真实 \texorpdfstring{$\zeta$}{zeta} 的可量化边界（用于对齐与路线设计）}
为了把“向 RH 推进”的外部进展接入本文的谱语言，这里列出若干在 2024--2025 仍可直接调用的硬输入（均为可核验的显式常数/阈值）：
\begin{itemize}
  \item \textbf{验证高度（partial RH）：}已严格验证到高度 \(3\cdot 10^{12}\)，即所有满足 \(0<\Im(\rho)\le 3\cdot 10^{12}\) 的非平凡零点都在临界线 \(\Re(\rho)=\tfrac12\) 上（\cite{PlattTrudgian2020RH3e12}）。
  \item \textbf{de Bruijn--Newman 常数：}Rodgers--Tao 证明 \(\Lambda\ge 0\)（\cite{RodgersTao2018NewmanNonnegative}），而 RH 等价于 \(\Lambda\le 0\)，因此 RH 等价于 \(\Lambda=0\)。同时 Polymath15 给出无条件上界 \(\Lambda\le 0.22\)，并给出依赖进一步数值验证的条件改进（\cite{Polymath2019NewmanUpper}）。
  \item \textbf{临界线上零点比例：}已有结果表明超过 \(5/12\) 的零点位于临界线（\cite{PrattRoblesZaharescuZeindler2018FiveTwelfths}）。该类“比例推进”可被视为把偏离临界线的零点集合压稀的一条可累积路线。
  \item \textbf{显式零点计数误差：}Bellotti--Wong 改进了 \(N(T)\) 的显式误差界，并给出具体常数（\cite{BellottiWong2024ImprovedNT}），可直接用于 Turing 方法/数值核验与显式应用。
  \item \textbf{Turing 方法关键量的改进：}Amberger 给出 \(|S_1(t_2)-S_1(t_1)|\) 的改进界，从而改进 Turing 方法所需的块长度阈值（\cite{Amberger2025S1}）。
  \item \textbf{\(\Re(s)\approx 1\) 的显式控制：}Cully-Hugill--Leong 与 Leong 分别给出 \(\zeta'(s)/\zeta(s)\) 与 \(1/\zeta(s)\) 在 \(\sigma\) 接近 \(1\) 时的显式上界（\cite{CullyHugillLeong2023Zeta1Line,Leong2024LogDerivRecip}），可作为零自由区与误差项常数的工程输入。
\end{itemize}

\begin{definition}[核版 Newman 阈值（可检验参数）]\label{def:kernel-newman-threshold}
给定一个带参数的核族 \(M(u)\)（例如由势函数/温度化门 \(\mathrm{PROJ}_u\) 或角色扭曲生成；见定义 \ref{def:pom-proj-temp} 与 \ref{def:pom-proj-chi}），记主谱半径为 \(\lambda_1(u)\)，最坏非平凡扭曲谱半径为 \(\Lambda(u)\)。
定义核版 Newman 阈值为
\[
\beta_\ast:=\inf\{\beta\in\RR_{\ge 0}:\ \Lambda(e^{-\beta})\le \lambda_1(e^{-\beta})^{1/2}\}.
\]
当 \(\beta_\ast=0\) 时，原点切片（\(u=1\)）已满足核版 GRH 的平方根界；当 \(\beta_\ast>0\) 时，需向零温方向变形（\(\beta\) 增大、\(u=e^{-\beta}\) 变小）后才进入平方根区间。
\end{definition}

\begin{remark}[可复现输出：纯碰撞轴的核版 Newman 阈值（有限模数代理）]\label{rem:kernel-newman-threshold-repro}
以真实输入 40 状态核的三轴扭曲族为例，固定 \((q,r)=(1,1)\)，并在第三轴取纯碰撞可观测量 \(N_2\)（即 \(\xi=\mathbf{1}_{d=2}\)），对每个奇素数 \(p\) 考察复合扭曲
\(
u=\exp(-\beta)\,\omega_p^k
\)
（\(k\not\equiv 0\!\!\pmod p\)）。定义
\[
\lambda_1(\beta):=\rho(M(\exp(-\beta))),\qquad
\Lambda_p(\beta):=\max_{k\neq 0}\rho\bigl(M(\exp(-\beta)\omega_p^k)\bigr),
\]
并令 \(\beta_\ast(p)\) 为满足 \(\Lambda_p(\beta)\le \lambda_1(\beta)^{1/2}\) 的最小 \(\beta\)（若存在）。
表 \ref{tab:real_input_40_kernel_newman_threshold} 给出若干 \(p\) 的数值扫描结果：若在扫描区间内未观测到满足平方根界的 \(\beta\)，则给出下界 \(\beta_\ast(p)>\beta_{\max}\)，并同时列出扫描网格上 \(\Lambda_p/\sqrt{\lambda_1}\) 的最小值及其发生位置（脚本 \texttt{scripts/exp\_real\_input\_40\_kernel\_newman\_threshold.py} 一键生成）：
\input{sections/generated/tab_real_input_40_kernel_newman_threshold}
\end{remark}

\begin{remark}[从 toy-RH 到“角度熵”：一个可操作的升级方向]\label{rem:finite-rh-angle-entropy}
命题 \ref{prop:finite-rh-triangle-dict} 表明：有限维 toy-RH 的“临界线位置”只取决于归一化谱的模长约束（\( |\mu|\le \sqrt\lambda \)）。
若要在有限模型中模拟更接近经典 \(\zeta\) 的“零点密度/统计形态”，则必须让临界模长（\(|\mu|=\sqrt\lambda\)）上的特征值携带大量不同的角度，从而把单一晶格振荡推广为多频振荡叠加（推论 \ref{cor:finite-rh-explicit-formula} 的结构推广）。
在本文的可编译口径下，这一目标可通过增加 lift/角色轴束（例如 profinite 同余轴或更高维 abelian 扭曲族）来实现，并可用“单位圆谱角分布的熵/多样性”作为可复现的实验指纹。
\end{remark}

\begin{remark}[平方根界的可审计证书：Lyapunov/内积构造]\label{rem:finite-rh-lyapunov}
核版 (G)RH 的平方根界是一个谱半径不等式。除直接对角化/求特征值外，还可用“构造内积”的方式给出可审计证书：
对任意矩阵 $A$ 与阈值 $r$，$\rho(A)<r$ 等价于存在正定矩阵 $H\succ 0$ 使 $A^\ast H A\prec r^2 H$（命题 \ref{prop:lyapunov-spectral-radius-certificate}）。
因此在扭曲分量 $A=M_{K,\rho}(u)$ 上，若能显式给出证书矩阵 $H_{\rho,u}$，就可把平方根界写成一条可复核的不等式，而不依赖数值特征值排序的稳健性。
\end{remark}

\paragraph{primitive 轨道计数曲线与 \texorpdfstring{$x/\log x$}{x/log x} 的常数因子}
令
\[
\Pi(N):=\sum_{n\le N}p_n
\]
为按长度累计的 primitive 轨道数（按轨道计）。由 $p_n\sim \lambda^n/n$ 的“末项主导 + 几何尾”可得
\[
\Pi(N)\sim \frac{\lambda}{\lambda-1}\cdot \frac{\lambda^N}{N}.
\]
在本核上 $\lambda=\varphi^2$ 且 $\lambda-1=\varphi$，故
\[
\Pi(N)\sim \varphi\cdot\frac{\varphi^{2N}}{N}=\frac{\varphi^{2N+1}}{N}.
\]
若定义“轨道范数” $N(\tau):=\lambda^{|\tau|}$ 并令
\[
\pi(x):=\#\{\tau\ \mathrm{primitive}:\ N(\tau)\le x\},
\]
则 $\pi(\lambda^N)=\Pi(N)$，从而
\[
\pi(x)\sim \frac{\lambda\log\lambda}{\lambda-1}\cdot \frac{x}{\log x}.
\]
对 $\lambda=\varphi^2$，该常数因子化简为
\[
\frac{\lambda\log\lambda}{\lambda-1}
=\frac{\varphi^2\log(\varphi^2)}{\varphi}
=2\varphi\log\varphi.
\]

定义轨道 Mertens 和
\[
\mathcal{M}(N):=\sum_{n\le N}\frac{p_n}{\lambda^n}
=\log N+\mathsf{M}_{\varphi^2}+o(1),
\]
其中常数项可写为
\[
\mathsf{M}_{\varphi^2}
=\gamma+\sum_{n\ge 1}\left(\frac{p_n}{\lambda^n}-\frac{1}{n}\right),
\]
$\gamma$ 为 Euler--Mascheroni 常数。数值上
\[
\sum_{n\ge 1}\left(\frac{p_n}{\lambda^n}-\frac{1}{n}\right)\approx -0.254312661901716,
\qquad
\mathsf{M}_{\varphi^2}\approx 0.322903002999817.
\]
另外记 $P(z)=\sum_{n\ge 1}p_n z^n$，则在 $z\to z_\star=\lambda^{-1}$ 时
\[
P(z)=-\log(1-\lambda z)+\log\mathfrak{M}+o(1),
\]
其中
\[
\log\mathfrak{M}\approx -0.254312661901716,\qquad
\mathfrak{M}\approx 0.7754493106833835.
\]
由因式分解可得主极点 $z_\star=\lambda^{-1}=\varphi^{-2}$，并且
\[
\zeta_M(z)=\frac{C}{1-\lambda z}+O(1),\qquad
C=\lim_{z\to z_\star}(1-\lambda z)\,\zeta_M(z)=\frac{47+21\sqrt5}{100}\ (\approx 0.9395742753).
\]

\begin{proposition}[主极点留数常数的闭式（可审计）]\label{prop:real-input-40-residue-constant}
沿用上文 \(\lambda=\rho(M)=\varphi^2\) 与 \(z_\star=\lambda^{-1}\) 的记号，并定义主极点留数常数
\[
C_{\mathrm{real}}:=\lim_{z\to z_\star}(1-\lambda z)\,\zeta_M(z).
\]
则 \(C_{\mathrm{real}}\) 具有闭式
\[
\boxed{\;
C_{\mathrm{real}}=\frac{47+21\sqrt5}{100}=\frac{13+21\varphi}{50}\; }.
\]
\end{proposition}
\begin{proof}
由 \(\zeta_M(z)=1/\Delta(z)\) 且 \(z_\star\) 为 \(\Delta\) 的简单零点，得
\(
\zeta_M(z)=\frac{1}{-\Delta'(z_\star)}\cdot\frac{1}{z-z_\star}+O(1)
\)。
又 \(1-\lambda z=-\lambda(z-z_\star)+O((z-z_\star)^2)\)，故
\(
C_{\mathrm{real}}=\frac{\lambda}{\Delta'(z_\star)}=\frac{1}{-z_\star\Delta'(z_\star)}
\)。
将 \(\Delta(z)=\det(I-zM)\) 的因子分解代入并化简即可得到所示代数闭式。证毕。
\end{proof}
进一步由 Witt 分解/ Möbius--log 反演（与 \S\ref{subsec:zeta-online-kernel} 相同口径）可得一个直接由 $\zeta_M$ 主极点读出的 Abel/有限部闭式：
\[
\log\mathfrak{M}
=\log C+\sum_{m\ge 2}\frac{\mu(m)}{m}\log\zeta_M(z_\star^{\,m}),
\qquad z_\star=\lambda^{-1}.
\]
右端把有理函数 $\zeta_M$ 代入代数数 $z_\star^{\,m}$，由于 $z_\star<1$，该级数以指数速度收敛。

\paragraph{显式 Möbius--$\zeta$ 级数与对数乘积（仅用 \texorpdfstring{$\det(I-zM)$}{det} 的因式分解）}
由
\[
\zeta_M(z)=\frac{1}{\Delta(z)},\qquad
\Delta(z)=\det(I-zM)=(1-z)^2(1+z)^3(1-3z+z^2)(1+z-z^2),
\]
以及 $z_\star=\lambda^{-1}\in(0,1)$，可将上式等价写成只含 $\Delta$ 的形式
\[
\log\mathfrak{M}
\;=\;
\log C-\sum_{m\ge 2}\frac{\mu(m)}{m}\log \Delta(z_\star^{\,m}).
\]
进一步把 $\Delta$ 的因子逐项展开，可得一个完全“正实数对数”的闭式：
\[
\log\mathfrak{M}
=\log C
-\sum_{m\ge 2}\frac{\mu(m)}{m}\Bigl[
2\log(1-z_\star^{\,m})
3\log(1+z_\star^{\,m})
\log(1-3z_\star^{\,m}+z_\star^{\,2m})
\log(1+z_\star^{\,m}-z_\star^{\,2m})
\Bigr].
\]
该式在数值上极稳定：因为 $z_\star^{\,m}$ 指数趋零，且每项只需对正数取对数。
等价地，可写成 Euler--Möbius 乘积（将“有限部常数”完全压缩为一个可比标量）：
\[
\mathfrak{M}
=C\prod_{m\ge 2}\zeta_M(z_\star^{\,m})^{\mu(m)/m}
=C\prod_{m\ge 2}\Delta(z_\star^{\,m})^{-\mu(m)/m}.
\]
\begin{lemma}[莫比乌斯坍缩：平凡因子的显式求和]\label{lem:mobius-collapse}
对 $|z|<1$ 有恒等式
\[
\sum_{m\ge 1}\frac{\mu(m)}{m}\log(1-z^m)=-z,\qquad
\sum_{m\ge 1}\frac{\mu(m)}{m}\log(1+z^m)=z-z^2.
\]
因此
\[
\sum_{m\ge 1}\frac{\mu(m)}{m}\log\Bigl((1-z^m)^{-a}(1+z^m)^{-b}\Bigr)=(a-b)z+bz^2.
\]
\end{lemma}

\begin{remark}[剥离平凡因子后的非平凡部分]
将引理 \ref{lem:mobius-collapse} 作用于 $(1-z)^2(1+z)^3$，其 Möbius--$\log$ 求和贡献可精确压缩为
\[
(-z_\star+3z_\star^2),
\]
其中 $-z_\star$ 来自 $(1-z)^2$ 与 $(1+z)^3$ 的线性合并项，而 $+3z_\star^2$ 是 $(1+z)^3$ 在修正恒等式中的二次修正项。
因此“剥离平凡因子”不仅产生线性项，还会带来一个可显式计算的二次项；这不会改变主极点位置，但会在 Abel 有限部常数的压缩表达中产生可观测的常数级修正。
\end{remark}

\begin{proposition}[约化 Mertens 常数：去平凡因子后的唯一贡献]\label{prop:real-input-40-reduced-mertens}
设
\[
\Delta(z)=\det(I-zM)=\Delta_{\mathrm{triv}}(z)\,\Delta_{\mathrm{nt}}(z),
\qquad
\Delta_{\mathrm{triv}}(z)=(1-z)^a(1+z)^b,
\]
其中 $a,b\in\ZZ_{\ge 0}$。假设主极点 $z_\star\in(0,1)$ 来自非平凡块，即 $\Delta_{\mathrm{nt}}(z_\star)=0$ 且为简单零点，而 $\Delta_{\mathrm{triv}}(z_\star)\neq 0$。
令
\[
\log\mathfrak{M}_{\mathrm{nt}}
:=\log C_{\mathrm{nt}}-\sum_{m\ge 2}\frac{\mu(m)}{m}\log \Delta_{\mathrm{nt}}(z_\star^{\,m}),
\qquad
C_{\mathrm{nt}}:=\frac{1}{-z_\star\,\Delta_{\mathrm{nt}}'(z_\star)}.
\]
则 Abel 有限部常数满足严格分解
\[
\boxed{\;
\log\mathfrak{M}
\;=\;
\log\mathfrak{M}_{\mathrm{nt}}
\;+\;
(a-b)\,z_\star
\;+\;
b\,z_\star^2.
\;}
\]
特别地，除去一个完全显式的“平凡修正项”\((a-b)z_\star+bz_\star^2\) 后，常数项的真正机制信息只依赖于非平凡谱块 $\Delta_{\mathrm{nt}}$。
\end{proposition}
\begin{proof}
由 $\Delta=\Delta_{\mathrm{triv}}\Delta_{\mathrm{nt}}$ 且 $\Delta_{\mathrm{nt}}(z_\star)=0$ 为简单零点，可得
\(
\Delta'(z_\star)=\Delta_{\mathrm{triv}}(z_\star)\Delta_{\mathrm{nt}}'(z_\star)
\)，
从而主留数常数满足
\(
C=\frac{1}{-z_\star\Delta'(z_\star)}=C_{\mathrm{nt}}/\Delta_{\mathrm{triv}}(z_\star)
\)。
将
\[
\log\mathfrak{M}=\log C-\sum_{m\ge 2}\frac{\mu(m)}{m}\log\Delta(z_\star^{\,m})
\]
代入并分解出 $\Delta_{\mathrm{triv}}$ 与 $\Delta_{\mathrm{nt}}$，得到
\[
\log\mathfrak{M}
=\log\mathfrak{M}_{\mathrm{nt}}
-\log\Delta_{\mathrm{triv}}(z_\star)
-\sum_{m\ge 2}\frac{\mu(m)}{m}\log\Delta_{\mathrm{triv}}(z_\star^{\,m}).
\]
对 $\Delta_{\mathrm{triv}}(z)=(1-z)^a(1+z)^b$，用引理 \ref{lem:mobius-collapse} 的两条恒等式并注意到
\(
\sum_{m\ge 2}=\sum_{m\ge 1}-\text{(}m=1\text{项)}
\)，
可直接化简得到
\[
-\log\Delta_{\mathrm{triv}}(z)-\sum_{m\ge 2}\frac{\mu(m)}{m}\log\Delta_{\mathrm{triv}}(z^{m})
=(a-b)z+bz^2.
\]
取 $z=z_\star$ 即得结论。证毕。
\end{proof}

\begin{definition}[Artin 符号通道的 Mertens 常数]\label{def:real-input-40-artin-sign-mertens}
在张量拆分中，因子
\(
\Delta_{\mathrm{twist}}(z)=1+z-z^2
\)
等价于
\(
\det(I-z(-F))=\det(I+zF)
\)
（其中 \(F=\begin{pmatrix}1&1\\1&0\end{pmatrix}\)）。
记
\[
\zeta_{-F}(z):=\frac{1}{1+z-z^2},
\qquad
\lambda=\varphi^2,\quad z_\star=\lambda^{-1}.
\]
定义其 Abel/Mertens 有限部常数（纯 twist 因子）为
\[
\boxed{\;
K_{-F}
:=\sum_{m\ge 1}\frac{\mu(m)}{m}\log\zeta_{-F}(z_\star^{\,m})
=-\sum_{m\ge 1}\frac{\mu(m)}{m}\log\!\bigl(1+z_\star^{\,m}-z_\star^{\,2m}\bigr)\; }.
\]
\end{definition}

\begin{remark}[可复现数值：\texorpdfstring{$K_{-F}$}{K_-F}]\label{rem:real-input-40-artin-sign-mertens-value}
由表 \ref{tab:real_input_40_finite_part_split} 的拆分，\(\log\mathfrak{M}_{\mathrm{twist}}\) 正是上述 \(K_{-F}\)（它完全由 \(1+z-z^2\) 这一非主极点因子贡献）。
为便于复核，我们同时给出一个高精度常数表：
\input{sections/generated/tab_real_input_40_artin_sign_mertens}
\end{remark}

此外，由 $\sqrt5=2\varphi-1$，留数常数亦可写成
\[
C=\frac{47+21\sqrt5}{100}=\frac{13+21\varphi}{50}.
\]
